% -----------------------------------------------------------
\section{Comforter}
% -----------------------------------------------------------
	\label{COMFORTER}
	
% % ----------------------
% \subsection{See also}
% % ----------------------

% % ----------------------
% \subsection{1828 American Dictionary of the English Language} % *1828
% % ----------------------
	% \label{COMFORTER-1828}

	% \begin{compactitem}
		% \item
	% \end{compactitem}

% % ----------------------
% \subsection{Oxford English Dictionary} % *OED
% % ----------------------
	% \label{COMFORTER-OED}

	% \begin{compactitem}
		% \item[]
	% \end{compactitem}
	
% % ----------------------
% \subsection{Restoration Lexicon} % *RL *RESTORATION LEXICON
% % ----------------------
	% \label{COMFORTER-OED}

	% \begin{compactitem}
		% \item[]
	% \end{compactitem}

% ----------------------
\subsection{Scriptural Usage} % *SCRIPTURES
% ----------------------
	\label{COMFORTER-SCRIPTURES}

	% -----
	\subsubsection{Old Testament} % *OT
	% -----
		\label{COMFORTER-OT}
	
		% --
		\paragraph{KJV}
		% --
			\label{COMFORTER-OT-KJV}
	
			\begin{tabular}{ll}
				Total occurrences in the OT & 9
			\end{tabular}
			
			\begin{compactdesc}
				\item[{\hyperref[2SAMUEL-10-3]{2 Samuel 10:3}}]
				\item[{\hyperref[1CHRONICLES-19-3]{1 Chronicles 19:3}}]
				\item[{\hyperref[JOB-16-2]{Job 16:2}}]
				\item[{\hyperref[PSALMS-69-20]{Psalms 69:20}}]
				\item[{\hyperref[ECCLESIASTES-4-1]{Ecclesiastes 4:1}}] \textbf{(2)}
				\item[{\hyperref[LAMENTATIONS-1-9]{Lamentations 1:9}}]
				\item[{\hyperref[NAHUM-3-7]{Nahum 3:7}}]
			\end{compactdesc}
			
		% % --
		% \paragraph{Hebrew Bible}
		% % --
			% \label{COMFORTER-HEBREW}
		
			% %
			% \subparagraph{\texorpdfstring{\he{sample word}}{}} % paste actual hebrew text here
			% %
				% \label{PAGA} % whatever is in step bible without periods
			
				% \begin{compactdesc}
					% \item[HALOT , PDF ] \textbf{} % Use the 1994 PDF
						% \begin{compactitem}
							% \item[1]
						% \end{compactitem}
					% \item[BDB , PDF ] \textbf{}
						% \begin{compactitem}
							% \item[1]
						% \end{compactitem}
					% \item[DCH PDF ] \textbf{} % don't include the actual page number, they repeat from volume to volume
						% \begin{compactitem}
							% \item[1]
						% \end{compactitem}
				% \end{compactdesc}
				
				% \begin{compactdesc}
					% \item[Some OT uses] \textbf{}
						% \begin{compactdesc}
							% \item[]
						% \end{compactdesc}
				% \end{compactdesc}
			
		% % --
		% \paragraph{Septuagint}
		% % --
			% \label{COMFORTER-LXX}
		
			% %
			% \subparagraph{\texorpdfstring{\gr{sample word}}{}}
			% %
		
	% -----
	\subsubsection{New Testament} % *NT
	% -----
		\label{COMFORTER-NT}
	
		% --
		\paragraph{KJV}
		% --
			\label{COMFORTER-NT-KJV}		
	
			\begin{tabular}{ll}
				Total occurrences in the NT & 4
			\end{tabular}
			
			\begin{compactdesc}
				\item[{\hyperref[JOHN-14-16]{John 14:16}}] \phantom{.}
					\begin{compactitem}
						\item See also verses 17--18 here.
					\end{compactitem}
				\item[{\hyperref[JOHN-14-26]{John 14:26}}]
				\item[{\hyperref[JOHN-15-26]{John 15:26}}]
				\item[{\hyperref[JOHN-16-7]{John 16:7--11}}]
			\end{compactdesc}
			
		% --
		\paragraph{Greek New Testament} % *GREEK
		% --
			\label{COMFORTER-GREEK}
		
			%
			\subparagraph{\texorpdfstring{\gr{παράκλητος}}{}}
			%
				\label{PARAKLETOS}
				
				\begin{tabular}{ll}
					Total occurrences in the NT & 5
				\end{tabular}
				
				\begin{compactdesc}
					\item[All NT uses] \phantom{.}
						\begin{compactdesc}
							\item[{\hyperref[JOHN-14-16]{John 14:16}}]
							\item[{\hyperref[JOHN-14-26]{John 14:26}}]
							\item[{\hyperref[JOHN-15-16]{John 15:16}}]
							\item[{\hyperref[JOHN-16-7]{John 16:7}}]
							\item[{\hyperref[1JOHN-2-1]{1 John 2:1}}]
						\end{compactdesc}
				\end{compactdesc}
			
				\begin{compactdesc}
				
					\item[BDAG 680b] originally meant in the passive sense...(but \gr{παράκλητος} is restored from \gr{παρακλος}, and the restoration is uncertain), `one who is called to someone's aid'. Accordingly Latin writers commonly rendered it, in its NT occurrences, with `advocatus'...but the technical meaning `lawyer', `attorney' is rare...In the few places where the word is found in pre-Christian and extra-Christian literature as well it has for the most part a more general sense: \textbf{one who appears in another's behalf, \textit{mediator, intercessor, helper}}...the Greek interpreters of John's gospel understood it in the active sense...in our literature the active sense \textit{helper, intercessor} is suitable in all occurrences of the word...in \hyperref[1JOHN-2-1]{1 John 2:1}...Christ is designated as \gr{παράκλητος}: \gr{παράκλητον ἔχομεν πρὸς τὸν πατέρα Ἰησοῦν Χριστὸν δίκαιον} \textit{we have Jesus Christ the righteous one, who intercedes for us}. The same title is implied for Christ by the \gr{ἄλλος παράκλητος} of \hyperref[JOHN-14-16]{John 14:16}. It is only the Holy Spirit that is expressly called \gr{παράκλητος} = \textit{Helper} in the Fourth Gospel.
						% (Note that this entry doesn't use the 1, 1a format of other entries.)
					
					
					
					\item[CGL 1069b, PDF 1109] \textbf{}
						\begin{compactitem}
							\item[1] one who is called in (to help), \textbf{supporter} (in a lawcourt) (Demosthenes)
							\item[2] (referring to the Holy Spirit) \textbf{counsellor, advocate} (New Testament)
						\end{compactitem}
					\item[LSJ 1313a, PDF 1384] \textbf{}
						\begin{compactitem}
							\item[I.1] \textit{called to one's aid}, in a court of justice: as substantive, \textit{legal assistant, advocate}
							\item[I.2] \textit{summoned}, \gr{δοῦλοι}
							\item[II] \textit{intercessor}; hence in NT, \gr{Παράκλητος}, of the Holy Spirit (\hyperref[JOHN-14-16]{John 14:16}, cf. \hyperref[1JOHN-2-1]{1 John 2:1})
						\end{compactitem}
					\item[Other sources] \phantom{.}
						\begin{compactitem}
							\item For what might be the most helpful discussion of the meaning of this word, see Brown's Appendix V of AB 29a, starting at PDF 604; this link goes to an annotated excerpt with all the pages you need: \href{https://drive.google.com/file/d/10zIKL33glyccw-vZwPUUZJatlqmsOOdh/view?usp=sharing}{here}.
							\item From Smalley's 1978 book \textit{John: Evangelist and Interpreter}, \href{https://drive.google.com/file/d/10_iMGwsHXcFinOCoO7j2GGI6oFtGnCa9/view?usp=sharing}{an excerpt on the ``paraklete.''}
						\end{compactitem}
				\end{compactdesc}
		
	% -----
	\subsubsection{Book of Mormon} % *BOM
	% -----
		\label{COMFORTER-BOM}
	
		\begin{tabular}{ll}
			Total occurrences in the BOM & 1
		\end{tabular}
		
		\begin{compactdesc}
			\item[{\hyperref[MORONI-8-26]{Moroni 8:26}}]
		\end{compactdesc}
		
	% -----
	\subsubsection{Doctrine and Covenants} % *DC
	% -----
		\label{COMFORTER-DC}
	
		\begin{tabular}{ll}
			Total occurrences in the DC & 23
		\end{tabular}
		
		\begin{compactdesc}
			\item[{\hyperref[DC21-9]{DC 21:9}}] \phantom{.}
				\begin{compactitem}
					\item The JSP version is slightly different.
				\end{compactitem}
			\item[{\hyperref[DC24-5]{DC 24:5}}] 
			\item[{\hyperref[DC28-1]{DC 28:1}}] 
			\item[{\hyperref[DC28-4]{DC 28:4}}] 
			\item[{\hyperref[DC31-11]{DC 31:11}}] 
			\item[{\hyperref[DC35-19]{DC 35:19}}]
			\item[{\hyperref[DC36-2]{DC 36:2}}]
			\item[{\hyperref[DC39-6]{DC 39:6}}]
			\item[{\hyperref[DC42-16]{DC 42:16}}]
			\item[{\hyperref[DC42-17]{DC 42:17}}]
			\item[{\hyperref[DC47-4]{DC 47:4}}]
			\item[{\hyperref[DC50-14]{DC 50:14}}]
			\item[{\hyperref[DC50-17]{DC 50:17}}]
			\item[{\hyperref[DC52-9]{DC 52:9}}]
			\item[{\hyperref[DC75-10]{DC 75:10}}]
			\item[{\hyperref[DC75-27]{DC 75:27}}]
			\item[{\hyperref[DC79-2]{DC 79:2}}]
			\item[{\hyperref[DC88-3]{DC 88:3}}] \textbf{(2)}
			\item[{\hyperref[DC88-4]{DC 88:4}}]
			\item[{\hyperref[DC90-11]{DC 90:11}}]
			\item[{\hyperref[DC90-14]{DC 90:14}}]
			\item[{\hyperref[DC124-97]{DC 124:97}}]
		\end{compactdesc}
		
	% -----
	\subsubsection{Pearl of Great Price} % *PGP
	% -----
		\label{COMFORTER-PGP}
	
		\begin{tabular}{ll}
			Total occurrences in the PGP & 1
		\end{tabular}
		
		\begin{compactdesc}
			\item[{\hyperref[MOSES-6-61]{Moses 6:61}}]
		\end{compactdesc}
		
% % ----------------------
% \subsection{Key Phrases} % *KEYPHRASES *PHRASES
% % ----------------------
	% \label{COMFORTER-PHRASES}

	% \begin{compactdesc}
		% \item[] \textbf{\#times} | list
			% % \begin{compactdesc}
				% % \item[Bible anchor verses] \phantom{.}
					% % \begin{compactdesc}
						% % \item[Romans 5:5] \gr{}
					% % \end{compactdesc}
				% % \item[Commentary] ---
			% % \end{compactdesc}
	% \end{compactdesc}
	
% % ----------------------
% \subsection{Bible Dictionaries} % *DICTIONARIES
% % ----------------------
	% \label{COMFORTER-BD}

% % ----------------------
% \subsection{Restoration Usage} % *RESTORATION
% % ----------------------
	% \label{COMFORTER-RESTORATION}

	% % -----
	% \subsubsection{Joseph Smith} % *JS
	% % -----
		% \label{COMFORTER-JS}
	
		% \begin{compactdesc}
			% \item[{\hyperref[KEY]{Date/Source}}]
		% \end{compactdesc}
		
	% % -----
	% \subsubsection{Temple Ordinances}
	% % -----
		% \label{COMFORTER-TEMPLE}
	
		% \begin{compactdesc}
			% \item[{\hyperref[KEY]{Date/Source}}]
		% \end{compactdesc}
	
	% % -----
	% \subsubsection{Brigham Young} % *BY
	% % -----
		% \label{COMFORTER-BY}
	
		% \begin{compactdesc}
			% \item[{\hyperref[KEY]{Date/Source}}]
		% \end{compactdesc}
	
% ----------------------
\subsection{Commentary and Studies} % *COMMENTARY *STUDIES
% ----------------------
	\label{COMFORTER-COMMENTARY}

	% % -----
	% \subsubsection{Key Points}
	% % -----
		% \label{COMFORTER-KEYS}
	
		% \begin{compactitem}
			% \item
		% \end{compactitem}
		
	% -----
	\subsubsection{The functions of the ``Comforter'' in restoration scripture}
	% -----
		\label{COMFORTER-functions}
		
		Here is a list of the functions or actions attributed to the ``Comforter'' in various places in the scriptures. Some of these could probably be collapsed into each other but for the sake of completeness I tried to separate them out as much as possible.
		
		\begin{compactitem}
			\item To be the means of giving words to people
				\begin{compactitem}
					\item DC 21:9; note that this happens ``through me,'' meaning ``through Christ''
					\item DC 24:5; they're ``things'' here, and not ``words,'' and they're given for the purpose of writing them
					\item DC 47:4 also involves writing
					\item DC 52:9
					\item DC 124:97
				\end{compactitem}
			\item To manifest that Jesus Christ was crucified for sins
				\begin{compactitem}
					\item DC 21:9
					\item Also DC 42:17, though indirectly
					\item DC 90:11, sort of
				\end{compactitem}
			\item To be the means of teaching people 
				\begin{compactitem}
					\item DC 28:1; it could be that the Comforter is delivering words to Oliver to use in teaching, or that it's the means by which the teaching is occurring, or that it's delivering the words to the listeners, like at 2 Nephi 33:1
					\item DC 36:2, to teach the ``peaceable things of the kingdom''; also DC 39:6,
					\item DC 50:14, to teach ``the truth''; also DC 79:2 and DC 124:97
					\item Possibly DC 52:9
					\item DC 75:10
					\item DC 90:11
				\end{compactitem}
			\item To lead people to speak or teach in meeting settings
				\begin{compactitem}
					\item DC 28:4; maybe also DC 42:16
				\end{compactitem}
			\item To tell people what to do and where to go
				\begin{compactitem}
					\item DC 31:11
					\item DC 75:27
					\item DC 79:2
				\end{compactitem}
			\item To reveal the fulness of the gospel and/or the mystery of things that have been sealed.
				\begin{compactitem}
					\item DC 35:19; note that this verse reads, ``Wherefore, watch over him that his faith fail not, and it shall be given by the Comforter, the Holy Ghost, that knoweth all things.'' The pronoun referent for ``it'' is not clear---it could be the ``faith'' spoken of in this verse, but it seems more likely that it's ``the fulness of my gospel'' in verse 17 or the ``mystery'' in verse 18. The other early sources don't seem to shed a lot of light on this.
					\item DC 90:14
				\end{compactitem}
			\item The means of baptism by fire
				\begin{compactitem}
					\item DC 39:6
				\end{compactitem}
			\item To show all things
				\begin{compactitem}
					\item DC 39:6; similar to 2 Nephi 32
				\end{compactitem}
			\item To know all things
				\begin{compactitem}
					\item DC 35:19, DC 42:17
				\end{compactitem} 
			\item To bear record of the Father and of the Son
				\begin{compactitem}
					\item DC 42:17
					\item DC 90:11
				\end{compactitem}
			\item To be the means for preaching the gospel
				\begin{compactitem}
					\item DC 50:14, DC 50:17
				\end{compactitem}
			\item To abide in our hearts
				\begin{compactitem}
					\item DC 88:3
				\end{compactitem}
			\item To be a sign of the promise God gives of eternal life
				\begin{compactitem}
					\item DC 88:4
				\end{compactitem}
		\end{compactitem}